\section{Analytic entry with one magnetic insertion}
\label{sec:entry}

The preparation in Section~\ref{sec:preparation} expresses the observed
height interaction as an exact sum of local patterns. We now estimate
that sum in the full analytic norm of \eqref{inp:norm}. The distinction
between an ordinary pattern and a pattern meeting the magnetic source
is important: ordinary activities will be arbitrarily small, whereas
the single activity carrying the source will have a polynomial bound
in the observation scale. No smallness of that insertion is required.

We use the notation of Section~\ref{sec:preparation}. In particular,
$s$ is the observation-cell side, $R$ is the tile parameter,
$P=R^5$, $m=qRs$, and $J_G^l$ is the cutoff interaction associated
with the affine observation energy $E_G^l$. The polynomial covariance
below $m$, its continuation on zero modes, and the analytic norm are
those fixed in Section~\ref{sec:inputs}. The common constant of the
reference field is uniform modulo $2\pi/\alpha$.

The regulator must be chosen to accommodate the large-observation
saving. By \eqref{prep:face-saving} and
\eqref{eq:affine-shape-bound}, there is a fixed $h_b\in(0,1)$ such
that
\[
 \frac{J_G^l(y)}{Z_G^0(0)}
 \le \exp\left\{\frac{h_b}{2}E_G^l(y)
       +C\log(2s)\,
            \mathbf 1_{\{\text{pierced star in }G\}}\right\}.
\]
The constants here are independent of $R,P,s$. Fix
$h_{\rm ent}$ with $h_b<h_{\rm ent}<1$. An
\emph{entry-admissible reserve} will mean a choice
\begin{equation}\label{eq:entry-admissible-reserve}
 h_{\rm ent}<h_1<h_*<1/p_2,
                    \qquad 1<p_1<p_2.
\end{equation}
Such choices exist by taking $p_1,p_2$ sufficiently close to $1$.
We make this choice, with room for any fixed finite collection of
stronger regulator regimes, before choosing $L$ or activity weights.
The proof below keeps all actual quadratic charges below
$h_{\rm ent}$.

In a pierced square, the \emph{compulsory cover} is the set of entry
blocks covering the mark, a fixed $CP$-cell neighborhood of it,
and the prescribed collars. Here $C$ is a sufficiently large fixed
geometric constant. The cover is declared occupied whether or not
an isolation error occurs there, and is also called the \emph{token}.
Its size can depend on $R$ but not on $s$ or on the square side.

\begin{proposition}[Mesoscopic analytic entry]\label{prop:entry}
Choose the regulator reserves as in
\eqref{eq:entry-admissible-reserve}, and fix compatible parameters
for the local maps of~\cite{B}, including the block factor $L$, the
support geometry, the activity weight $\mathfrak A$, and the analytic
radius. A fixed finite collection of stronger weights and
entry-admissible regulator regimes may be specified. For every
$\epsilon>0$, one can choose a fixed dyadic
$q$, an arbitrarily large dyadic $R$, and then $s_0$, such that the
following conclusions hold for every compatible dyadic $s\ge s_0$,
with $P=R^5$ and $m=qRs$.

\begin{enumerate}
\item In the ordinary case, averaging the observation noise and the
reference covariance below $m$ gives local activities $K_m$ satisfying
\begin{equation}\label{eq:entry-small}
                 \norm{K_m}_{m,\mathfrak A,V}\le\epsilon.
\end{equation}
This holds on finite plane supports, on compatible tori, and on free
squares whose sides are divisible by $sP$ and are at least $sP$.
Up to a positive field-independent scalar, the corresponding subset
partition function is exactly the averaged cutoff interaction.
Activities with their required clearance from a wall or a torus
identification are the corresponding plane activities. The construction
is real, even, and covariant under the block translations and square
symmetries that preserve the chosen hierarchy.

\item In a pierced free square in which the compulsory footprint and
its required collars fit, the same scalar as in its ordinary
counterpart may be extracted. The remaining partition function is
\begin{equation}\label{eq:entry-linear-partition}
                    \mathcal Z_m^{(1)}(K_m^l,I_m;\psi),
\end{equation}
with exactly one insertion activity $I_m$. Its supports contain the
mark block and the compulsory cover just specified. The background
activities $K_m^l$ are arbitrarily small in their prescribed norms,
and
\begin{equation}\label{eq:entry-insertion}
                    \norm{I_m}_{m,\mathfrak A,V}
                         \le C_R s^{C_R}.
\end{equation}
Here and below $C_R$ may also depend on the fixed downstream norm and
geometric parameters, but not on $s$ or on the square side. Ordinary
occupied supports remain separated from the insertion by the original
compatibility rule.

\item A support is called \emph{shift-good} if its spanning rectangle,
enlarged by $Hm$, misses the magnetic mark. The fixed constant $H$
is chosen from the local-map geometry. On such supports the pierced
background is the ordinary background evaluated at the local lift
$\psi+\lambda$, where
\[
                    \partial\lambda=\eta_l/\alpha.
\]
On interior supports near the mark, including insertion supports, the
activities are copies of the pierced-plane calculation evaluated at
$\psi+\gamma_M$, with $\gamma_M$ and its derivative bounds as in
\eqref{eq:lift-difference} and \eqref{eq:lift-derivatives}. Copying always
includes the collars and numerical dependence required by the activity.
\end{enumerate}
The assertions hold simultaneously in the fixed finite collection of
reserve regimes. Their constants are uniform after the geometric
parameters have been fixed.
\end{proposition}

We prove the proposition by adapting the real-pattern and analytic
estimates of~\cite[Section 5, Lemma
5.4]{X}, together with the normalization and connected
pattern arguments of~\cite[Section 9]{B}. The point of
the proof is to check the hypotheses of those arguments for the
unstrengthened XY law with cutoffs and for its affine magnetic sector.
The Gaussian coefficient in every comparison below is $\alpha^2=8\pi$.

\subsection{Bounded observations and conditional restoration}

An isolation error has two positive terms: the restoration kernel
before a cut, and the restoration kernel after the cut multiplied by
its Gaussian proxy. Denote these kernels by $K_o$ and $K_c$, and the
proxy by $B_i$. After conditioning on the guard heights and prescribing
the observations, its relative absolute error is
\begin{equation}\label{eq:entry-relative-error}
                  \delta_i=
            \frac{|K_o-B_iK_c|}{K_o+B_iK_c}\le1.
\end{equation}
The kernels in this formula include the face cutoffs retained in their
respective graphs. A cutoff common to both terms cancels from the
ratio. In particular, inserting these cutoffs does not change the
conditional sampling of guard heights at fixed observations.

First suppose that the observations range over a fixed compact set
and that the patch has no magnetic source. On this set
$\chi_{\sqrt{s}}(d_f)$ tends uniformly to $1$, with positive values;
its logarithm therefore contributes a uniformly vanishing error to
\eqref{eq:entry-relative-error}. Other positive real factors only
decrease when multiplied by cutoffs. For a selected isolated tile,
whose factor is its normalized interaction minus $1$, the same compact
convergence applies to that difference.

The remaining convergence inputs are precisely the mixed-smear,
separated-pin, and conditional-pin limits of~\cite{X}, with the
observation mass retained. Core pins and surrounding guards have fixed
positive clearance in tile units. All cuts follow whole observation
cells, and free-wall patches are clipped with the same actual wall.
These are the mixed polygonal geometries of the cited limits. The
observation mass controls the constant on each component. Compactness
of observation data is understood modulo component periods, using the
joint components of all faces appearing in a restoration test.

For clarity, the conditional-pin input is convergence in probability
under the relevant guard laws, not convergence uniformly at arbitrary
guard values. In the proof of~\cite[Equation
(5.20)]{X}, conditional terms are integrated under
either of the two restoration topologies. The restoration comparisons
of~\cite[Lemma 3.9]{X} give exactly these changes
of sampling law. The local-limit inputs required there are smear
limits with observation mass, separated hard-pin and conditional-pin
limits, their versions with fixed observation tilts, and the
corresponding finite-geometry Gaussian statements. For the present
unstrengthened law these are supplied directly by
\cite[Theorem 3.1 and the following restoration
estimates]{X}; the additional averaged-bond precision used in that
paper's strengthened-law entry lemma is absent here.

These statements are uniform on the compact observation sets just
specified. One way to see the needed uniformity is to let the
observation data themselves converge in such a compact set. Sampling
tilts then vary continuously, with uniform integrability supplied by
the observation exponential-moment bounds. For each finite list of
conditional Laplace terms, differences at converging arguments tend
to zero in probability: their derivatives are bounded by conditional
moments on a fixed larger compact set. The finite Gaussian terms and
proxies converge in the same way, and the limiting positive
normalizers stay bounded away from zero. Compactness now gives the
claim.

Here is the quantitative estimate supplied by these limits. Let
$\mathcal E_{\rm box}$ be the sum of squared affine face differences
on the underlying portion of a $P$-box in the padded pattern.
It may include faces from a patch whose difference is being tested.
Recenter one observation in each joint component into $[-\pi,\pi]$
by a height period. Paths within the box then give
\[
       \norm y_2\le C P^C
                          (1+\sqrt{\mathcal E_{\rm box}(y)}).
\]
Thus a fixed bound on $\mathcal E_{\rm box}$ gives a compact
observation set. Call an error \emph{deep} if its core, open patch,
and testing guard lie inside the bounding box with the prescribed
positive clearances. At fixed $R,P$, there are finitely many box
geometries, positive topologies, and choices of deep errors, and their
number $j_{\rm deep}$ is at most a fixed polynomial $N_P$ in $P$.

For a deep error $i$, write
$\mathcal D_i=\log(K_o/(B_iK_c))$.
Decompose this logarithm into its centered zero-guard value, the
change of centered restoration normalization with guard data, and
the two conditional observation Laplace terms minus the proxy
energy change. The first term is $O(e^{-cR})$ after refinement.
The second lies between $\log r_{H_i}$ and zero, where $r_{H_i}$
is the conditional full-core pin ratio on a neighborhood $H_i$
containing the restored bonds. These are the restoration and
pin-interaction estimates just cited.

In the limiting Gaussian problem, let $u_{i,c}$ be the cut harmonic
prediction from the guard data. Set
$\varepsilon_i=\vartheta_i u_{i,c}$, where $\vartheta_i$ equals
one near $H_i$ and vanishes before the guard. The norm
$\norm{\cdot}_{\mathcal T}$ is the continuum
gradient-plus-unit-cell-observation energy norm on the cut patch,
with gradient coefficient $\alpha^{-2}$, corresponding to
\eqref{prep:massive-form}. The guarded localization estimate
of~\cite[Equation (9.23)]{B} gives, for every $r\ge2$,
\begin{equation}\label{eq:entry-prediction}
 \left(\E_{\tau,y}^{\rm Gauss}
                  \norm{\varepsilon_i}_{\mathcal T}^{\,r}\right)^{1/r}
 \le C\sqrt r\,P^C e^{-cR}
                        (1+\sqrt{\mathcal E_{\rm box}(y)}).
\end{equation}
Here $\tau$ may be either restoration topology, and the expectation
is over its Gaussian guard law with fixed observation tilt $y$.
The constants and polynomial degree are independent of $r$.
Indeed the centered prediction is $\mathsf A g$ in the energy
Hilbert space, where $g$ has identity covariance and
$\norm{\mathsf A}_{\rm HS}\le CP^C e^{-cR}$.
Diagonalization and Gaussian moments bound its $L^r$ energy norm
by $C\sqrt r\,\norm{\mathsf A}_{\rm HS}$. Restoring bonds decreases
the guard covariance. The observation tilt adds a deterministic
prediction bounded by $CP^Ce^{-cR}\norm y_2$, proving
\eqref{eq:entry-prediction} under both laws.

Subtracting $\varepsilon_i$ from a cut minimizer makes it admissible
for restoration without changing the guard data. Orthogonal
projection in the energy space gives
\[
       -2\log r_{H_i}\le\norm{\varepsilon_i}_{\mathcal T}^2
\]
and bounds the linear conditional-Laplace discrepancy by
$C\norm y_2\norm{\varepsilon_i}_{\mathcal T}$. The conditional
quadratic discrepancy from the open-patch proxy is at most
$Ce^{-cR}\norm y_2^2$, by the same guarded localization.
If restoration joins free faces, first prescribe finitely many
smooth guard probes: the cut and restored minimizers belong to
nested admissible energy spaces, and the correction makes the
cut minimizer trace-matching. Increasing these finite probe lists
gives the assertion. Thus the calculation uses finite-energy
predictions, not pointwise values of a rough continuum trace.

Since $\delta_i\le\min(1,|\mathcal D_i|)$, these estimates give
Gaussian moment bounds of the form
\[
 \norm{\delta_i}_{L^r(\tau,y)}
       \le e^{-c_0R}P^C(1+r)^C
                              (1+\mathcal E_{\rm box}(y))^C.
\]
The microscopic conditional approximation transfers these bounded
moments under either topology. For this transfer use the normalized
conditional restoration ratio: the conditional expectation of the
positive restoration factor in the centered cut law, divided by
its value at zero guard. This ratio lies between $r_{H_i}$ and $1$,
and the limiting Gaussian $r_{H_i}$ is strictly positive.
Exceptional sets therefore have vanishing
probability under both tilted laws, and there $\delta_i\le1$
suffices. The cutoff ratios add only uniformly vanishing logarithms
on the compact observation set.

Let $\mu_{\tau,y}$ be a positive topology's conditional guard law
in the box at fixed observations, and let $\mathcal I_{\rm deep}$
be its retained deep errors. H\"older's inequality, without an
independence assumption on these errors, now yields
\begin{equation}\label{eq:entry-good-box}
 \E_{\mu_{\tau,y}}\!
       \left[\prod_{i\in\mathcal I_{\rm deep}}\delta_i\right]
       \le C\exp\left\{-cRj_{\rm deep}
                          +R^{-12}\mathcal E_{\rm box}(y)\right\}.
\end{equation}
For $j_{\rm deep}\ge1$, use moment order
$\max(2,j_{\rm deep})$. The polynomial loss is at most
\[
       [P^C(1+j_{\rm deep})^C
                   (1+\mathcal E_{\rm box})^C]^{j_{\rm deep}}.
\]
Maximizing its energy-dependent part after multiplication by
$e^{-R^{-12}\mathcal E_{\rm box}}$ costs
$e^{O(j_{\rm deep}\log R)}$, since
$j_{\rm deep}\le N_P=\operatorname{poly}(P)$ and $P=R^5$.
Large $R$ absorbs this loss in the exponential rarity. For
$j_{\rm deep}=0$ the left-hand side is one. The remaining
quadratic allowance is precisely
\begin{equation}\label{eq:entry-good-allowance}
                         R^{-12}\mathcal E_{\rm box}.
\end{equation}

This transfer needs no convergence rate uniform in $R$. Fix $R,P$
and the real-data threshold first, and prove the Gaussian version
of \eqref{eq:entry-good-box} with a stronger exponential constant.
Compactness and the finite list of geometries make the refinement
errors, including exceptional-set probabilities, smaller than any
prescribed positive number. Choose this number below the finite
minimum of the remaining slack for $1\le j_{\rm deep}\le N_P$,
then take $s$ sufficiently large. The prediction constants and
exponential gain have already been fixed; only the refinement
threshold depends on this transfer tolerance. For an isolated
selected tile the normalized Gaussian interaction is one, so the
same compact convergence makes its error from one at most
$e^{-cR}$. These are the event and tile gains used below in the
real pattern estimate.

If a simply connected observation box misses the source, the ray
connection is removed by an integral gauge constant on each
observation cell. The preceding argument then applies without change,
with $d_f^2$ in the bounding face cost. Enlarged guards and counting
halos can be increased by fixed factors. At an actual free wall,
including a corner, the guards are clipped with that wall; no periodic
replacement of the boundary condition is used.

\subsection{Centered constants and cuts into bounded boxes}

The estimates must retain the centered normalization even in the
magnetic calculation. In each pattern telescope only the true centered
partition ratios in its affected halos. These are $Z_G^0(0)$ without
cutoffs. The restoration and pin-interaction estimates of~\cite{X}
therefore give the same bounded exponential cost per marked event as in
\cite[Equations (5.16) and
(5.17)]{X}. A fixed distinguished halo may have an extra
$e^{C_R}$ cost. Nothing is telescoped through an unaffected portion of
the volume.

After this accounting, each positive topology has the real field
factor
\begin{equation}\label{eq:entry-topology-factor}
       \frac{J_G^l(y)}{Z_G^0(0)}
       =e^{E_G^l(y)/2}
          \prod_{f\in\mathcal F(G)}\chi_{\sqrt{s}}(d_f)
                         \frac{Z_G^l(y)}{Z_G^0(0)},
\end{equation}
multiplied by a conditional expectation of products of
\eqref{eq:entry-relative-error} and by the small energy-error factors.
Here $\mathcal F(G)$ is the set of observation faces retained in that topology.
The double telescoping of normalization constants in
\cite[Lemma 9.3]{B} now uses the affine matrices
of Section~\ref{sec:preparation}. Their local curl terms cancel
exactly: the two paths compare the same deletions with the same
nearby affine data and connection. At artificial boundaries the two
paths are clipped identically in the unchanged collar. Only the
exponentially small coefficients on face squares remain.

To estimate a topology, cut it into observation boxes of side $P$.
The shifted comparison \cite[Equation (9.20)]{B}
applies with the connection and with prescribed guard values by
Proposition~\ref{inp:comparisons}. Its guard probabilities still use centered
partition sums. Drop the cutoffs on the new box cuts. The conditional
kernels at deep errors, and hence their cutoff ratios, are unchanged.
The Gaussian localization estimate gives, after averaging translates
of the cuts, the logarithmic cost
\begin{equation}\label{eq:entry-seam}
                         C(R/P)\mathcal E,
\end{equation}
where $\mathcal E$ is a face-square cost confined to the original
padded support. Indeed each face is within the fixed tile-size
neighborhood of a seam for a proportion $O(R/P)$ of the translates;
the exponentially decaying matrix tails have the same bound after
summation. A local curl term can add $C\log(2s)$ only in the box
containing the magnetic junction. No other elementary cell junction
has nonzero circulation.

The same averaging retains a fixed positive proportion, for example
one half, of the deep errors while satisfying
\eqref{eq:entry-seam}, after increasing its constant. One can first
separate the common upper charge for the magnetic curl and then apply
the ordinary averaging argument. This is the only exceptional seam
charge associated with the source.

\subsection{The compulsory cover and large observations}

In the pierced calculation use the compulsory cover defined before
Proposition~\ref{prop:entry}. Choose its constant $C$ large enough that
every bounding box belonging to a component disjoint from this cover
can be tested away from the source. There is exactly one token in a pattern.

In the token component use the good-box estimates outside the fixed
number of boxes meeting the source. In those exceptional boxes take
only $\delta_i\le1$. They contain at most $\operatorname{poly}(P)$
events, and the missing gains can therefore be paid by a constant
depending on $R$. We next give a bound that applies both to these
boxes and to boxes with large observations.

The cutoff large-field estimate and \eqref{eq:affine-shape-bound}
give the bound used in choosing $h_b$ above:
\begin{equation}\label{eq:entry-broad-bound}
 \frac{J_G^l(y)}{Z_G^0(0)}
 \le \exp\left\{
       \frac{h_b}{2}E_G^l(y)
             +C\log(2s)\,\mathbf 1_{\{\text{pierced star in }G\}}
                    \right\}.
\end{equation}
To check the strict inequality $h_b<1$, combine
\eqref{eq:large-observation} with the Gaussian decay of the cutoff
on the remaining faces. The logarithm of the right-hand side of
\eqref{eq:entry-topology-factor} is at most
\[
              \tfrac12 E_G^l(y)-c\sum_{f\in\mathcal F(G)}d_f^2.
\]
Since \eqref{eq:affine-shape-bound} bounds $E_G^l$ by a fixed multiple
of that face sum plus the displayed curl charge, a fixed fraction
of $E_G^l/2$ is saved. This fraction is independent of $s,R,P$.
The initial choice of $h_b$ can be made positive by enlarging it
while preserving $h_b<1$. For a selected isolated-tile difference the same argument
allows the sum of this bound and $1$.

It remains to express these affine costs in the homogeneous energy
used in the norm. Write
\[
                         x=y-B_sr,
\]
taking $r=0$ in the ordinary sector, and let $D$ be the underlying
padded support. Restrictions of a trial field for $E_D(x)$ show that
the sum of the cut homogeneous observation energies is at most
$E_D(x)$. A quadratic inequality in \eqref{eq:shape-energy}, followed
by \eqref{eq:affine-shape-bound}, bounds the affine part from $B_sr$
by $O(\operatorname{poly}(P))$ per broad box, and by
$O_R(\log(2s))$ for the token calculation. Here the affine face
differences of $B_sr$ are uniformly bounded by the source estimates
of Section~\ref{sec:preparation}. Choose the quadratic-inequality
parameter so that its increase in $h_b$ is less than one quarter
of $h_{\rm ent}-h_b$. This is a fixed choice made before $R$;
its deterministic cost only changes the constants depending on $R$.

Small face costs are handled by the two-cell Poincare estimate. Their
deterministic affine parts cost their coefficient times the number
of faces. In particular, \eqref{eq:entry-seam} adds
$O((R/P)R^2)$ per affected halo, and
\eqref{eq:entry-good-allowance} adds $O(R^{-12}P^2)$ per occupied
bounding box. Link costs have their exponentially small summable
coefficients. Both large-observation thresholds and the regulators
are evaluated on real data. If a box is declared large because
$\mathcal E_{\rm box}>U_{\rm big}^2$, insert
\begin{equation}\label{eq:entry-bad-payment}
             1\le
       \exp\{\epsilon_1(\mathcal E_{\rm box}-U_{\rm big}^2)\},
\end{equation}
with fixed sufficiently small $\epsilon_1>0$. Choose it so that its
increase in the energy coefficient uses less than another quarter
of $h_{\rm ent}-h_b$. Its positive part
uses only a further small increase of the energy coefficient and
$O(\operatorname{poly}(P))$ per such box. Its negative part will
pay the lost good-box gains.

The remaining seam, good-box, and link coefficients tend to zero
as $R$ increases. Choose $R$ so that their sum uses less than the
remaining half of $h_{\rm ent}-h_b$. The actual energy coefficient
$h_{\rm used}$ therefore satisfies
\[
        h_{\rm used}<h_{\rm ent}<h_1<h_*<1/p_2,
                   \qquad 1<p_1<p_2,
\]
with the regulator parameters already fixed in
\eqref{eq:entry-admissible-reserve}. This ordering is uniform for
large $s$ and applies in each of the finitely many reserve regimes.

\subsection{Gaussian averaging and the real pattern bound}

We now average the observation noise and the reference modes below
$m$. This step uses the contraction of the homogeneous observation
energy and the localized determinant estimate in
\cite[Lemma 9.2]{B}. It is important to apply the
trace estimate to the actual quadratic costs before replacing them
by their common energy majorant. A broad box has only $O(P^2)$
observation coordinates, so its contribution has rank $O(P^2)$.
This remains true for a box meeting the magnetic source. Every
Gaussian adjacent-cell difference has bounded variance: use the
two-cell flow from Section~\ref{sec:preparation}, now in the
ordinary Gaussian field. Hence the remaining trace estimate is
exactly the localized trace estimate of~\cite[Equation
(5.24)]{X}. Links retain their exponential path costs.

After fixing norm parameters and then large $R,P$, choose the finite
$U_{\rm big}$ so that the negative term in
\eqref{eq:entry-bad-payment} pays the broad-box trace, all polynomial
box costs, and all lost item gains. Only after this choice do we
invoke the compact-data estimates and increase $s$. The identities
$P=R^5$ give
\[
       (R/P)R^2=R^{-2},\qquad R^{-12}P^2=R^{-2},
\]
so the seam and good-box allowances leave strict room for the
exponential gains.

For a pattern $\pi$, let $j_1$ count its isolation errors and $j_2$
its selected tiles. Let $d_\ell$ denote the cell-distance parameter
of a link, including its numerical dependence radius. If $X$ is the
covered block support and $F_\pi$ its averaged field factor, the
resulting real estimate has the form
\begin{align}\label{eq:entry-real-pattern}
 |F_\pi(\psi)|
 &\le C^{j_1+j_2}e^{-cR(j_1+j_2)}
           \prod_\ell Ce^{-c(R+d_\ell)}
                     W_m^{\kappa_0}(X,\psi)e^{V_0(X,\psi)},
       \qquad \kappa_0<\kappa.
\end{align}
For a pattern with the token, multiply the right-hand side by
$C_Rs^{C_R}$. The fixed number of exceptional boxes and the
$O_R(\log(2s))$ curl charge are precisely the sources of this
factor. There is only one such factor. The sum over choices of
small- and large-observation boxes costs a fixed exponential per
ordinary item, or a fixed $C_R$ in the token component.

If a favorable cut translate was selected according to the real
argument, dominate that selection by the sum over all translates.
The number of choices gives at most a polynomial-in-$R$ factor per
error, absorbed in the exponential item gain. No optimized-cut
comparison is needed for a pattern with no error event; fixed
token costs are already permitted. In particular, this step never
introduces a determinant whose rank grows with the microscopic
cell size $s$.

The same regulator argument applies on free squares. Restricting
to an open support lowers the homogeneous minimum energy. The
positive-order covariance estimates for shells and for the terminal
field follow from the binary-polynomial kernels of~\cite{B}, with
the convention of Section~\ref{sec:inputs} and even reflection at
free walls. The independent-normal factor from other covariance
splits is absent here. Enlarge supports only by aligned whole
cells. This proves the energy contraction, localized trace, and
finite-range additivity needed in the free-square estimates.

The strict reserve also allows real translations of the background.
If $f$ has scaled analytic test norm at most $1$, replacing $\psi$
by $\psi+xf$ in \eqref{eq:entry-real-pattern} costs at most
\begin{equation}\label{eq:entry-real-translation}
                             e^{Cx^2|X|},
\end{equation}
while leaving the final regulator evaluated at $\psi$. This is
the quadratic completion with reserve in
\cite[Equation (9.28)]{B}.

\subsection{Complex directions and analytic derivatives}

The real bound alone does not imply the full analytic norm. We obtain
a separate estimate on a fixed complex strip and then interpolate.
Let $z$ be a real observation direction. On the extended Gaussian
chain lattice, an affine connection and prescribed real observations
become a real linear tilt and a scalar after choosing coordinates
on the underlying centered lattice. The theta parity identity
of~\cite[Equation (9.29)]{B} therefore gives
\begin{equation}\label{eq:entry-complex-comparison}
           \frac{|Z_G^l(y+iz)|}{Z_G^l(y)}
                    \le\frac{Z_G^0(iz)}{Z_G^0(0)}.
\end{equation}
The fixed-graph Gaussian-chain limit preserves this inequality.
The right-hand side is the ordinary centered imaginary ratio for
the topology, without cutoffs. Its characteristic factor is positive,
and centered pinning increases that factor.

Discard the rarity of the errors, expand into positive topologies,
and apply \eqref{eq:entry-complex-comparison} to each whole integrated
topology. Use the argument of~\cite{X} from its complex-comparison
Equation (5.26) through its
strip estimate (5.27): pin whole observation-cell
strips of width $R/4$ along translated $P$-box walls. The mixed
limits on the resulting fixed pinned boxes give their Gaussian
imaginary ratios uniformly on each fixed compact set of imaginary
directions. The limiting ratios are strictly positive. The additional
imaginary Gaussian wall cost is
\begin{equation}\label{eq:entry-imaginary-wall}
          O_{H_{\rm strip}}(R/P)
                    \#\{\text{active observation cells}\}
\end{equation}
on a strip of fixed height $H_{\rm strip}$. We distinguish this
strip height from the geometric clearance $H$ in
Proposition~\ref{prop:entry}. The homogeneous imaginary quadratics
cancel against the same proxy quadratics as in the ordinary argument;
connection offsets are real.

The cutoff factors need a simpler estimate. For real $a,v$, their
defining Gaussian integral gives
\begin{equation}\label{eq:entry-complex-cutoff}
                 |\chi_{\sqrt{s}}(a+iv)|
                     \le e^{v^2/2}\chi_{\sqrt{s}}(a).
\end{equation}
If the field direction has bounded scaled analytic test norm, the
imaginary part of an adjacent observation difference is
$O_{H_{\rm strip}}(s/m)$. A halo has $O(R^2)$ cells, so the sum of
these extra logarithmic costs is $O_{H_{\rm strip}}(1)$ per halo.
It does not require a new wall estimate. The real parts of the
positive-topology estimates still have the bounds already proved:
on compact-data boxes use the same comparisons with rarity discarded,
and on large-data boxes retain their Gaussian completion and payments.
For a selected tile discard the subtraction from $1$ in this rough
bound. The token component may again pay its single factor
$C_Rs^{C_R}$.

Thus, for real $x,t$ with $|t|\le H_{\rm strip}$ and a direction
$f$ of scaled point norm at most $1$,
\begin{align}\label{eq:entry-rough-strip}
 |F_\pi(\psi+(x+it)f)|
 &\le
 \exp\{o(R)(j_1+j_2)+C_{H_{\rm strip}}(1+x^2)|X|\}
      \prod_\ell Ce^{-c(R+d_\ell)}
                      W_m^\kappa(X,\psi)e^{V_0(X,\psi)}.
\end{align}
A token multiplies this by $C_Rs^{C_R}$. The term $o(R)$ includes
the wall cost: per halo \eqref{eq:entry-imaginary-wall} is
$O(R^3/P)=O(R^{-2})$. Absolute Gaussian integration and analyticity
follow from the broad real estimate with reserve and this strip bound.

Choose $r_{\rm an}>4eh_0$, where $h_0$ is the analytic radius, and then
$H_{\rm strip}>4r_{\rm an}$, before choosing $R$. Apply the three-lines
theorem on both half-strips to
\[
                  F_\pi(\psi+zf)e^{-D_{\rm damp}z^2|X|},
\]
where the constant $D_{\rm damp}>C_{H_{\rm strip}}$ is used only in this
interpolation. The real boundary has
\eqref{eq:entry-real-pattern} and
\eqref{eq:entry-real-translation}; the outer boundary has
\eqref{eq:entry-rough-strip}. On $|z|\le r_{\rm an}$ the harmonic weight
of the real boundary is at least $3/4$. The exponential gains
therefore survive, with a smaller positive constant. Removing the
damping costs $e^{C|X|}$, which is absorbed by the item and link
gains because
\[
              |X|\le C(j_1+j_2)
                    +C\sum_\ell(1+d_\ell/R)
\]
in the ordinary case, with an additional fixed $C_R$ for a token.
Finally Cauchy's estimate followed by real polarization bounds the
order-$r$ analytic derivative contribution by the same pattern
bound times
\[
                         (h_0/r_{\rm an})^r\frac{r^r}{r!}.
\]
The series in $r$ converges. We have proved
\eqref{eq:entry-real-pattern} in the full analytic point norm,
with weakened but still positive exponential gains.

\subsection{Summation, exact copying, and translated norms}

All bounding boxes above use only the underlying portions of the
original padded activity supports. They are devices for estimating
a pattern, not new activity domains. A cut crossing a topology
may charge underlying faces there, so the estimate also controls
differences on an error patch. All good-error patches tested after
cutting are deep in their bounding boxes. Choose cut translates
on the tile grid and imaginary pin strips on the $R/4$ grid;
these choices preserve the tile-junction geometry. A group meeting
a free wall uses its clipped events and patches, with the plane
ordering restricted to present tiles. Its size and counting bounds,
including at corners, are the same bounds used above.

The connected-pattern summation following
\cite[Equation (9.34)]{B} now applies. Exponential item gains pay
the activity weights, halo covers, and polynomial labels. A link
with dependence radius $D'$ has a cover of
$O(1+D'/R)$ blocks and decay exponential in $R+D'$, as proved in
Section~\ref{sec:preparation}; its gain therefore also pays its
cover and radius labels. On a torus the last telescoping term
uses a winding label with the same bound. The finite-range Gaussian
averaging factorizes on compatible components with their original
exclusions. Every ordinary connected component contains an item
or a link and hence retains a small factor. By increasing $R$,
their sum gives \eqref{eq:entry-small}.

The component containing the compulsory cover has the bound
\eqref{eq:entry-insertion}; all other components are ordinary
background activities. No additional optional token is summed.
This gives exactly the coefficient partition function
\eqref{eq:entry-linear-partition}. The normalization constants were
centered throughout, so the extracted scalar is the ordinary one
on the same square. The absolute estimates justify every exchange
of pattern summation and Gaussian integration. In particular these
are convergent exact representations, not identities of formal
expansions.

Plane objects in this construction are defined on finite supports.
The tile and link rules can respect all block translations and
square symmetries, using symmetric spanning paths when necessary.
Padding an event cover includes the numerical rules in a sufficiently
large fixed multiple of its tile-size halo. Link covers also record
their numerical dependence. Consequently a support with its stated
clearance has exactly the same local calculation in the plane and
in the finite graph.

For the magnetic copying statement, recall that the harmonic energy
scalar has already been removed in \eqref{eq:interaction}. The
remaining localized preparation, including all links, depends on a
harmonic solution only through the arguments of its functions.
The polynomial reference covariances themselves copy with clearance.
Replacing $r_\infty$ by $r$ therefore changes an interior argument
precisely by $\gamma_M$ from \eqref{eq:lift-difference}. On a
shift-good support an integral gauge removes the connection and
gives the ordinary activity at the corresponding local lift.

There is one harmless convention needed for background activities
at entry. A support incompatible with the fixed compulsory token
footprint cannot appear in the coefficient partition function.
Define the background on such a support to be its prescribed ordinary
translate if it is shift-good, and zero otherwise. This leaves the
represented partition function unchanged and preserves the
pierced-plane copying rule. It also permits choosing $H$ from the
fixed map geometry, independently of $P$. Use identical positions
in the initial hierarchy in every comparison, and require the
inner supports of the token construction to fit with clearance.
The square applications have $M/(sP)\longrightarrow\infty$, so
this requirement is satisfied.

Lastly we check the norm convention for lifted backgrounds, rather
than merely their pointwise copying. On a shift-good support and
its collars the lift satisfies
\[
            m|\nabla\lambda|+m^2|\nabla^2\lambda|\le C.
\]
The ordinary regulator can use a doubled exponent, at a cost
$C^{|X|}$ on support $X$, as in~\cite[Proof of
Proposition 5.7]{B}. For the variational regulator, increase
its parameter $h_*$ slightly. The optimized quadratic is bounded
by local gradient energy and increases at least proportionally
with $h_*$; applying the quadratic inequality to the translation
again costs only $C^{|X|}$. A stronger untwisted activity weight
pays both losses. These changes use the finitely many reserve
regimes fixed at the start. Each downstream contraction may be
performed in its own fixed regime, with the same regulator at
successive steps; no reserve is consumed anew at every iteration.

Choose $h_b,h_{\rm ent}$ from the uniform preparation estimates,
then the analytic translation radii and entry-admissible regulator
reserves before
$L$, take $L$ sufficiently large for the desired local-map
contractions, and then take the activity weights sufficiently
large. Once these downstream parameters are fixed, choose $R,P$,
then $U_{\rm big}$, and finally $s_0$ as above. This order proves
all assertions of Proposition~\ref{prop:entry} simultaneously.
\qed
